%=============================================================================
% DISCUSSION AND CONCLUSION
%=============================================================================
\section{Discussion}
\label{sec:discussion}

\textbf{Why \model{} works, and what it cannot do.} Weekly dengue counts are close to a
random walk, so the most informative feature is the last value. A direct head must
carry that value through the encoder; STGAT, A3TGCN and DCRNN lose it and underfit.
\model{} hands the level to the output, so the encoder learns only the departure, and
the gate keeps a weak departure from pulling the forecast away from persistence. This is
why the same head repairs three different architectures, attention, recurrent and
diffusion based, by 6--16 RMSE. It also explains its limit: the anchor removes failure,
but it cannot create information. When the next large move is not in the inputs, as in
the 2026 outbreak, every model, \model{} included, lags it, and on RMSE no learned model
beats persistence on later data.

\textbf{Where physics helps a graph network, and where it does not.} On this target the
SEIR model cannot be the forecaster: a pure SEIR decoder inherits the poor
predictability of its force of infection, and 14--16\% of targets are out of its reach.
Inside the gated head it looked useful, because it made every encoder usable, but the
ablation shows that the anchor and gate do that work, and on later data the simulator
ties its no-physics twin. Coupling districts inside the dynamics does not help either,
because transmission changes have almost no local spatial structure at district-week
resolution. The graph itself adds 0.06 RMSE. The prior that matters is the one the data
already state: next week looks like this week.

\textbf{What the evaluation taught us.} Each of the following would have produced a
publishable positive result had it gone unchecked: the benchmark array leaks future data
through its row order and its climate channels; our own first SEIR-GNN pipeline compared
one origin against a three-origin baseline; the gated SEIR head's repair looked like
physics and was the anchor; the validation-selected model $B$ is worse than persistence
on test; and the SEIR arm's development advantage over persistence disappears on
2024--2026 data. We therefore recommend that dengue forecasting papers report
persistence, a matched no-physics control and, where possible, a prospective test.

\textbf{Limitations.} One country and one disease at district-week resolution. Three
origins cannot support claims at the origin unit, and the nine-origin test spans were
examined in earlier experiments, so they are not fresh. The prospective test is amended
(\S\ref{sec:newweeks}) and uses one encoder family; it is a single 2.4-year period
dominated by one outbreak year. The graph is the benchmark's hand-built adjacency, and
mobility data might couple districts differently. The SEIR branch uses fixed stage
durations and one initial susceptible fraction for every district.

\section{Conclusion}
We proposed \model{}, a persistence-anchored gated head for spatio-temporal dengue
forecasting that plugs into any encoder at the cost of one parameter, together with a
corrected, leakage-free dataset extended to 2026 with newly parsed reports. \model{} turns
three published ST-GNNs that are worse than persistence into usable forecasters, and
with a strong encoder it beats persistence on every paired run in the development
study. An ablation shows that the gain comes from the anchor and the gate; an SEIR
simulator in the same head adds nothing, on development data or on later weeks. On
weeks from 2024 to 2026 no learned model beats persistence on RMSE, though \model{} is
better on MAE and SMAPE. For district-level dengue in Sri Lanka, the useful prior is
persistence, not transmission dynamics, and the next gains must come from inputs that
anticipate outbreaks rather than from the architecture. The dataset, the notebook that
produces every development number, the pre-registrations and the experiment log are
available at \codeurl.
